% This file recreates the shared player-reference page. Its deliberately compact
% blocks are meant to sit in the center of the table during play, as in the source.
\begin{multicols}{2}
\DndrSectionBanner{Section 3: Comprehensive Player Reference Rules}
\textit{Keep this section accessible in the center of the table as a shared system reference guide during active gameplay.}

\subsection*{3.1 Standard Game Setup \& Campaign Scaling}
Decide how many Quests you would like to play, then build a Quest deck of that size.

\textbf{Short Adventure:} 2 Quests\\
\textbf{Medium Adventure:} 3 Quests\\
\textbf{Long Adventure:} 5 Quests

New groups should expect their first few games to take longer as they learn the rules and develop their storytelling style.

\subsection*{3.3 Reading Your Class Mat}
Your Class Mat is your character's home base. It tracks your unique Class Ability, Background, Applause to use, Signature Item, equipped Items, and health throughout the game.
\DndrGraphicPlaceholder{2.45in}

\subsection*{3.4 Applause Tokens}
Applause Tokens reward players for making the story more entertaining. Whenever someone makes the table laugh, tells a memorable flashback, or describes a particularly cool moment, toss them one of your Applause Tokens.

\textbf{Your own tokens are worthless to you.} To use your Class Ability, Signature Item, or an equipped Item, spend an Applause Token received from another player.

You may hold up to 5 Applause Tokens from other players. If a player at the limit would receive one, reward another player who helped make the moment memorable. When spent, return it to the player whose color it matches.

\subsection*{3.5 Gear \& Equipment Architecture}
All Items are activated the same way: spend an Applause Token, describe a brief flashback, and resolve the card's effects.

Most Items require at least one matching Skill Tag shown on your Background Card. If an Item has no Skill Tags, any character may use it.

\textbf{Signature Items} are equipped in the Signature Slot and represent iconic gear. They are never discarded after use and ignore all Skill Tag requirements.

Items come in three types:
\begin{itemize}
  \item \textbf{Basic Items (Blue):} Standard equipment and tools.
  \item \textbf{Rare Items (Red):} More powerful and unusual effects.
  \item \textbf{Instant Items (Green):} May be played during another player's or an enemy's turn and do not count against the normal item activation limit.
\end{itemize}
\end{multicols}
